% ============================================================
% MÓDULO: 5S, LEAN MANUFACTURING Y KAIZEN
% ============================================================
\clearpage
\section*{MÓDULO: 5S, LEAN MANUFACTURING Y KAIZEN}
\subsection*{Orden, eficiencia y mejora continua}

\begin{recuadro}
\textbf{Objetivo del módulo:} conocer las herramientas de gestión de origen japonés que usan
las industrias y los talleres para trabajar con orden, seguridad y menos desperdicio, y
aplicarlas en el pañol y el taller de la escuela.
\end{recuadro}

% ------------------------------------------------------------
\section*{1. LA METODOLOGÍA 5S}

Las \textbf{5S} son cinco palabras japonesas que comienzan con la letra ``S''. Forman un método
para lograr un lugar de trabajo \textbf{limpio, ordenado, seguro y productivo}. Es la base de
cualquier sistema de calidad: sin orden no hay mejora posible.

\begin{center}
\small
\begin{tabular}{|p{0.13\linewidth}|p{0.17\linewidth}|p{0.58\linewidth}|}
\hline
\textbf{Japonés} & \textbf{Significado} & \textbf{Qué se hace} \\ \hline
\textbf{Seiri} & Clasificar & Separar lo necesario de lo innecesario y retirar lo que no se
usa. Se usan \textbf{tarjetas rojas} para marcar los elementos dudosos. \\ \hline
\textbf{Seiton} & Ordenar & ``Un lugar para cada cosa y cada cosa en su lugar''. Lo más usado,
más cerca. Tableros de sombra (siluetas de herramientas), etiquetas, códigos de ubicación. \\ \hline
\textbf{Seiso} & Limpiar & Limpiar es también \textbf{inspeccionar}: al limpiar se descubren
pérdidas, fisuras, cables dañados, herramientas descalibradas. \\ \hline
\textbf{Seiketsu} & Estandarizar & Fijar por escrito cómo se mantienen las tres primeras S:
listas de verificación, fotos del ``estado correcto'', responsables y frecuencias. \\ \hline
\textbf{Shitsuke} & Disciplina & Convertir las 5S en hábito: auditorías periódicas, ejemplo de
los responsables, mejora del propio estándar. \\ \hline
\end{tabular}
\end{center}

\paragraph*{5S en la aviación.}
En un hangar las 5S no son solo estética: son \textbf{seguridad operacional}. El control de
herramientas evita que una llave olvidada dentro de un motor o de una bahía de aviónica
provoque un daño por objeto extraño (\textbf{FOD}, \textit{Foreign Object Damage}). Por eso
los talleres trabajan con tableros de sombra, herramientas numeradas, inventario al final de
cada tarea y registro de préstamos en el pañol.

% ------------------------------------------------------------
\section*{2. LEAN MANUFACTURING}

\textbf{Lean} (``esbelto'') es una filosofía de gestión surgida del \textbf{Sistema de
Producción Toyota}. Su idea central es \textbf{entregar valor al cliente eliminando todo lo que
no agrega valor}, es decir, el desperdicio (\textit{muda}).

\subsubsection*{2.1. Los 8 desperdicios}
\begin{enumerate}
  \item \textbf{Sobreproducción:} producir más o antes de lo necesario.
  \item \textbf{Esperas:} personas o máquinas paradas esperando material, información o una
        aprobación.
  \item \textbf{Transporte:} mover materiales de un lado a otro sin necesidad.
  \item \textbf{Sobreprocesamiento:} hacer más de lo que el cliente pide o necesita.
  \item \textbf{Inventario:} exceso de stock que inmoviliza dinero, ocupa lugar y se vence o
        se vuelve obsoleto.
  \item \textbf{Movimientos innecesarios:} buscar herramientas, caminar hasta el pañol varias
        veces.
  \item \textbf{Defectos:} errores que obligan a rehacer el trabajo.
  \item \textbf{Talento no aprovechado:} no escuchar las ideas de quienes hacen el trabajo.
\end{enumerate}

\subsubsection*{2.2. Herramientas Lean}
\begin{itemize}
  \item \textbf{Mapa de flujo de valor (VSM):} dibujo de todo el proceso, desde el pedido hasta
        la entrega, marcando qué pasos agregan valor y cuánto tiempo se pierde en esperas.
  \item \textbf{Justo a tiempo (JIT):} producir o comprar solo lo necesario, en la cantidad
        necesaria y en el momento necesario.
  \item \textbf{Kanban:} tarjetas o señales visuales que indican cuándo reponer un material. En
        el pañol, por ejemplo, cuando la caja de un componente llega a la marca de stock
        mínimo se dispara el pedido de compra.
  \item \textbf{Poka-yoke:} dispositivos ``a prueba de errores'' que impiden equivocarse, como
        conectores que solo entran en una posición.
  \item \textbf{Gestión visual:} tableros, colores y señales que muestran el estado del trabajo
        de un vistazo.
  \item \textbf{SMED:} reducir el tiempo de cambio de una preparación a otra.
  \item \textbf{TPM (Mantenimiento Productivo Total):} los operarios participan en el cuidado
        básico de sus equipos para evitar paradas.
\end{itemize}

% ------------------------------------------------------------
\section*{3. KAIZEN: LA MEJORA CONTINUA}

\textbf{Kaizen} significa ``cambio para mejor''. Propone \textbf{mejoras pequeñas, constantes y
con participación de todos}, en lugar de esperar grandes inversiones o grandes cambios.

\subsubsection*{3.1. El ciclo PDCA (Deming)}
\begin{enumerate}
  \item \textbf{Planificar (Plan):} identificar el problema, medir la situación actual, buscar
        la causa raíz y proponer una mejora.
  \item \textbf{Hacer (Do):} probar la mejora en pequeña escala.
  \item \textbf{Verificar (Check):} medir los resultados y compararlos con la situación inicial.
  \item \textbf{Actuar (Act):} si funcionó, estandarizarla; si no, ajustar y volver a empezar.
\end{enumerate}

\subsubsection*{3.2. Herramientas para encontrar la causa raíz}
\begin{itemize}
  \item \textbf{Los 5 porqués:} preguntar ``¿por qué?'' sucesivamente hasta llegar a la causa
        de fondo. Ejemplo: \textit{La placa se entregó tarde. ¿Por qué? Faltaba un integrado.
        ¿Por qué? No se repuso a tiempo. ¿Por qué? Nadie sabía que quedaba poco stock. ¿Por
        qué? No hay registro de consumos. $\Rightarrow$ Solución: registrar las salidas del
        pañol y fijar un stock mínimo.}
  \item \textbf{Diagrama de Ishikawa (espina de pescado):} ordena las causas posibles en
        categorías: Mano de obra, Métodos, Materiales, Máquinas, Medición y Medio ambiente.
  \item \textbf{Diagrama de Pareto:} el 80\% de los problemas suele venir del 20\% de las
        causas; se atacan primero las más frecuentes.
\end{itemize}

\begin{recuadro}
\textbf{Kaizen y seguridad operacional.} La mejora continua y el Sistema de Gestión de la
Seguridad Operacional (SMS) comparten la misma lógica: detectar, analizar, corregir y
verificar. Un reporte de un error o de un peligro no es una acusación, sino una oportunidad de
mejora.
\end{recuadro}

% ------------------------------------------------------------
\section*{4. ACTIVIDADES DEL MÓDULO}

\begin{enumerate}
  \item \textbf{Jornada 5S en el pañol o el taller:} relevamiento fotográfico del estado
        inicial, aplicación de tarjetas rojas, ordenamiento, armado de un tablero de sombra y
        foto del estado final. Redacción de la lista de verificación (Seiketsu).
  \item \textbf{Cacería de desperdicios:} observar durante una clase de taller y registrar
        ejemplos de cada uno de los 8 desperdicios.
  \item \textbf{Kaizen en equipo:} elegir un problema real del taller o del proyecto final,
        analizarlo con 5 porqués e Ishikawa, proponer una mejora y presentarla siguiendo el
        ciclo PDCA.
  \item \textbf{Kanban de componentes:} definir stock mínimo y punto de pedido para cinco
        componentes del pañol. Este trabajo se continúa en el proyecto integrador con Odoo.
\end{enumerate}
